\documentclass[11pt]{article}

\usepackage[margin=1in]{geometry}
\usepackage[utf8]{inputenc}
\usepackage[T1]{fontenc}
\usepackage{hyperref}
\usepackage{xcolor}
\usepackage{enumitem}
\usepackage{listings}
\usepackage{parskip}
\usepackage{titlesec}
\usepackage{booktabs}

\definecolor{odekared}{RGB}{178,49,40}
\definecolor{codebg}{RGB}{245,245,245}
\definecolor{codeframe}{RGB}{210,210,210}

\hypersetup{
    colorlinks=true,
    linkcolor=odekared,
    urlcolor=odekared,
    citecolor=odekared,
    pdftitle={Maintaining the ODEKA Website},
    pdfauthor={ODEKA A.S.B.L}
}

\titleformat{\section}{\normalfont\Large\bfseries\color{odekared}}{\thesection}{1em}{}
\titleformat{\subsection}{\normalfont\large\bfseries}{\thesubsection}{1em}{}
\titleformat{\part}[display]{\normalfont\huge\bfseries\color{odekared}\centering}{Part \thepart}{0.5em}{}

\lstset{
    basicstyle=\ttfamily\small,
    backgroundcolor=\color{codebg},
    frame=single,
    rulecolor=\color{codeframe},
    breaklines=true,
    breakatwhitespace=false,
    columns=fullflexible,
    keepspaces=true,
    xleftmargin=8pt,
    xrightmargin=8pt,
    framextopmargin=6pt,
    framexbottommargin=6pt,
    showstringspaces=false,
}

\newcommand{\term}[1]{\textbf{#1}}
\newcommand{\file}[1]{\texttt{#1}}
\newcommand{\cmd}[1]{\texttt{#1}}

\title{\color{odekared}\Huge Maintaining the ODEKA Website\\[0.3em]
\normalsize\color{black} A Guide for Research Associates}
\author{ODEKA A.S.B.L}
\date{\today}

\begin{document}

\maketitle
\tableofcontents

\part{Getting Started with Quarto}

\section{Introduction}

This guide is for research associates and other ODEKA staff who need to update
the ODEKA website (\url{https://odeka-asbl.github.io}) --- adding a news post,
a job listing, a new team member, or a new project, for example.

The website is built with a tool called \term{Quarto}. You do not need to know
how to code to use it: pages are written mostly in plain text (with a little bit
of HTML sprinkled in for layout), and Quarto turns that text into the actual
website.

This first part of the guide covers two things:

\begin{enumerate}
    \item Setting up your computer with the tools you need (Section~\ref{sec:setup}).
    \item How Quarto works in general --- the concepts you need to understand
    before touching this specific website (Section~\ref{sec:quarto-basics}).
\end{enumerate}

Part II covers the specifics of \emph{this} website: its folder
structure, its custom design system, and step-by-step instructions for
common tasks (adding a news post, adding a team member, etc.).

\section{Setting Up Your Computer}
\label{sec:setup}

You only need to do this once. If a tool is already installed on your
computer, you can skip that step.

\subsection{Install Visual Studio Code}

Visual Studio Code (\term{VS Code}) is the text editor we use to write and
edit the website's files.

\begin{enumerate}
    \item Go to \url{https://code.visualstudio.com}.
    \item Download the version for your operating system (Mac, Windows, or
    Linux) and run the installer.
    \item Open VS Code once to confirm it launches correctly.
\end{enumerate}

\subsection{Install Quarto}

Quarto is the actual engine that converts our source files into the website.

\begin{enumerate}
    \item Go to \url{https://quarto.org/docs/get-started/}.
    \item Download and run the installer for your operating system.
    \item To confirm it installed correctly, open a terminal (on Mac:
    Terminal app; on Windows: Command Prompt or PowerShell) and type:
\begin{lstlisting}
quarto check
\end{lstlisting}
    This should print a summary confirming Quarto is installed, along with
    the version number, without any errors.
\end{enumerate}

\subsection{Install the Quarto Extension for VS Code}

This extension gives you syntax highlighting, a ``Preview'' button, and other
conveniences directly inside VS Code.

\begin{enumerate}
    \item Open VS Code.
    \item Click the Extensions icon in the left-hand sidebar (it looks like
    four small squares), or press \cmd{Cmd+Shift+X} (Mac) / \cmd{Ctrl+Shift+X}
    (Windows).
    \item Search for \term{Quarto}.
    \item Install the extension published by \term{Quarto}.
\end{enumerate}

Once installed, opening any \file{.qmd} file in VS Code will show a
``Preview'' button (a small icon with a magnifying glass, usually in the
top-right corner of the editor) that opens a live preview of that page.

\subsection{Install Git and Get Access to the Repository}

The website's files live in a \term{Git repository} hosted on
\term{GitHub}, at a URL your supervisor will share with you. Git is the
tool that tracks changes to the files and lets you publish (``push'') your
edits so they show up on the live website.

\begin{enumerate}
    \item Create a free account at \url{https://github.com} if you don't
    already have one, and ask your supervisor to give your account access
    to the repository.
    \item Install Git:
    \begin{itemize}
        \item \term{Mac}: open Terminal and type \cmd{git --version}. If
        Git isn't already installed, macOS will prompt you to install the
        Xcode Command Line Tools, which include it.
        \item \term{Windows}: download and install Git from
        \url{https://git-scm.com/downloads}.
    \end{itemize}
    \item (Optional but recommended for beginners) Install
    \term{GitHub Desktop} (\url{https://desktop.github.com}), a graphical
    app that lets you commit and push changes without typing Git commands
    directly. VS Code also has built-in Git support (the ``Source Control''
    icon in the sidebar) if you prefer to stay in one app.
\end{enumerate}

\subsection{Connecting VS Code to Your GitHub Account}
\label{sec:git-auth}

Before you can push changes, Git needs to know it's really you. Note that
GitHub no longer lets you sign in for this using your regular account
password --- you'll instead use one of the two methods below.

\paragraph{Option A: Sign in through VS Code (recommended).} This is the
easiest method and involves no typing of tokens at all.

\begin{enumerate}
    \item In VS Code, click the \textbf{Source Control} icon in the
    left-hand sidebar (it looks like a branching line).
    \item Click \textbf{Sign in} (or, if you don't see it there yet, wait
    until you perform your first push --- VS Code will prompt you
    automatically at that point).
    \item Your web browser will open a GitHub page asking you to log in
    (username and password there, on GitHub's own site) and then
    \term{authorize Visual Studio Code} to access your account.
    \item Click \textbf{Authorize}. You'll be sent back to VS Code, which
    is now connected to your GitHub account. You will not be asked again
    on this computer.
\end{enumerate}

If you installed GitHub Desktop instead, the equivalent step is:
\textbf{GitHub Desktop} $\to$ \textbf{Settings} (Mac) or \textbf{Options}
(Windows) $\to$ \textbf{Accounts} $\to$ \textbf{Sign in to GitHub.com},
which opens the same kind of browser sign-in window.

\paragraph{Option B: Personal Access Token (if you work from the terminal).}
If you type Git commands directly in a terminal rather than using VS
Code's Source Control panel or GitHub Desktop, the first time you
\cmd{git push} you'll be prompted for a username and password. Enter your
GitHub username as usual, but for the password, you must use a
\term{Personal Access Token} (a long, randomly generated code) instead of
your real password --- typing your real password here will simply fail.

To create one:

\begin{enumerate}
    \item Go to \url{https://github.com/settings/tokens}.
    \item Click \textbf{Generate new token} $\to$ \textbf{Generate new
    token (classic)}.
    \item Give it a name (e.g. ``ODEKA laptop''), set an expiration date,
    and check the box next to \textbf{repo}.
    \item Click \textbf{Generate token} at the bottom of the page, then
    \textbf{copy the token immediately} --- GitHub will only show it to
    you once.
    \item The next time Git asks for a password, paste this token instead.
    Your computer will normally remember it after that, so you won't need
    to repeat this often.
\end{enumerate}

Unless you have a specific reason to work from the terminal, Option A is
simpler and is what we recommend.

\subsection{Clone the Repository}

``Cloning'' means downloading a working copy of the website's files onto
your computer.

\begin{enumerate}
    \item Open a terminal and navigate to the folder where you want to keep
    the project, e.g.:
\begin{lstlisting}
cd Documents
\end{lstlisting}
    \item Clone the repository (replace the URL with the one your
    supervisor gives you):
\begin{lstlisting}
git clone https://github.com/odeka-asbl/odeka-asbl.github.io.git
\end{lstlisting}
    \item Open the resulting folder in VS Code: \textbf{File} $\to$
    \textbf{Open Folder...} and select it.
\end{enumerate}

You now have everything you need on your machine. The remainder of this
document explains how Quarto itself works.

\section{Quarto: The General Principles}
\label{sec:quarto-basics}

\subsection{What Quarto Is}

Quarto is a \term{static site generator}. That means: you write content in
simple source files, and Quarto ``renders'' (builds) those files into
plain HTML pages --- the actual files a web browser displays. Nothing runs
on a server when someone visits the site; every page is just a
pre-built HTML file sitting in a folder, served as-is. This makes the site
fast, simple to host (we use GitHub Pages, which is free), and hard to
break in unpredictable ways.

The two building blocks you'll work with constantly are:

\begin{itemize}
    \item \term{\file{.qmd} files} --- these are the source files for each
    page. ``qmd'' stands for Quarto Markdown. You write and edit these.
    \item \term{The \file{docs/} folder} --- this is where Quarto puts the
    finished HTML files after rendering. This is what actually gets
    published to the live website. \textbf{You should never hand-edit
    files inside \file{docs/}} --- they get overwritten every time the site
    is rendered again. Always edit the \file{.qmd} source files instead.
\end{itemize}

\subsection{Markdown, Briefly}

Most of the text content on the site is written in \term{Markdown}, a
simple way to format text using plain characters instead of a toolbar.
The basics:

\begin{lstlisting}
# This is a big heading
## This is a smaller heading

Regular paragraph text goes here. You can make text **bold**
or *italic*. You can add a [link](https://example.com), or an
image: ![alt text](path/to/image.jpg)

- A bullet point
- Another bullet point
\end{lstlisting}

Our website also mixes in small pieces of raw HTML directly inside the
\file{.qmd} files (for things like banners, custom layouts, and buttons)
--- this is normal and something we'll cover in Part II. You don't need to
know HTML deeply, but you will encounter it.

\subsection{The YAML Frontmatter}

Every \file{.qmd} file starts with a block between two lines of three
dashes (\cmd{---}). This is called the \term{frontmatter}, written in a
format called \term{YAML}, and it controls settings for that specific
page: its title, whether it has a table of contents, and so on.

\begin{lstlisting}
---
title: "About us"
toc: false
---

The rest of the page's content starts here, after the
closing `---`.
\end{lstlisting}

\subsection{The Project-Wide Configuration File}

At the root of the project sits a single file, \file{\_quarto.yml}, which
controls settings for the \emph{entire} website rather than one page. This
is where the navigation bar, the site title, and the overall output
folder are defined. For example, a simplified version of ours looks like:

\begin{lstlisting}
project:
  type: website
  output-dir: docs

website:
  title: "ODEKA"
  navbar:
    left:
      - href: 01_about/about.qmd
        text: About us
      - href: 02_people/people.qmd
        text: People

format:
  html:
    theme: cosmo
    css: styles.css
\end{lstlisting}

Notice how each navbar entry simply points to a \file{.qmd} file and gives
it a label. Adding a new top-level page to the site generally starts with
adding an entry here.

\subsection{Rendering the Site}

``Rendering'' is the act of converting your \file{.qmd} source files into
the finished HTML in \file{docs/}. You do this from the terminal, from
the project's root folder:

\begin{lstlisting}
quarto render
\end{lstlisting}

This rebuilds the \emph{entire} site. If you only changed one page and
want a faster rebuild, you can render just that file:

\begin{lstlisting}
quarto render 04_news/index.qmd
\end{lstlisting}

You need to render (and then commit and push, see Section~\ref{sec:publish})
before your changes actually appear on the live website.

\subsection{Previewing Your Changes Live}

Rendering the whole site every time you make a small edit is slow. Instead,
while you're actively working on a page, use \term{preview mode}, which
opens a local copy of the site in your browser and automatically refreshes
it every time you save a file:

\begin{lstlisting}
quarto preview
\end{lstlisting}

This will print a local address (something like
\url{http://localhost:4848}) and normally opens it in your browser
automatically. Leave this command running in the terminal while you work;
press \cmd{Ctrl+C} in the terminal to stop it when you're done.

If you installed the Quarto VS Code extension mentioned earlier, you can
also just click the \textbf{Preview} button at the top of an open
\file{.qmd} file, which does the same thing without needing the terminal.

\subsection{Publishing Your Changes}
\label{sec:publish}

The live website is served directly from the \file{docs/} folder in the
GitHub repository. Publishing a change means three steps:

\begin{enumerate}
    \item \term{Render} the site (or the page you changed), so
    \file{docs/} is up to date:
\begin{lstlisting}
quarto render
\end{lstlisting}
    \item \term{Commit} your changes --- this saves a labeled snapshot of
    what changed, both in the source files and in \file{docs/}:
\begin{lstlisting}
git add .
git commit -m "Add news post about the new field team"
\end{lstlisting}
    \item \term{Push} your commit to GitHub, which is what actually makes
    it live (usually within a minute or two):
\begin{lstlisting}
git push
\end{lstlisting}
\end{enumerate}

If you're using GitHub Desktop or VS Code's Source Control panel instead
of typing Git commands, these same three steps exist as buttons: review
your changed files, write a short commit message, click \term{Commit},
then click \term{Push}.

\subsection{How a Quarto \emph{Website} Project Is Organized}

A single Quarto document is just one file. A Quarto \emph{website} project
--- which is what we have --- is a folder containing many \file{.qmd}
files (one per page), plus:

\begin{itemize}
    \item \term{Shared styling}: a single \file{.css} file (ours is
    \file{styles.css}) that controls fonts, colors, spacing, and layout
    across every page, so the whole site looks consistent.
    \item \term{Listings}: some pages (like News or Jobs) don't list their
    content by hand. Instead, they automatically generate a list of posts
    from a folder, sorted by date, based on settings in that page's
    frontmatter. Adding a new post is often just a matter of adding a new
    file to the right folder --- the listing picks it up automatically.
    \item \term{Includes}: small snippets of HTML or JavaScript that get
    inserted into every page automatically (for example, code that
    powers a hover effect in our navigation bar). These live in their own
    files and are referenced from \file{\_quarto.yml}.
    \item \term{Images and other assets}: stored alongside the pages that
    use them, referenced by file path from within the \file{.qmd} files.
\end{itemize}

\subsection{A Typical Editing Workflow, Start to Finish}

Putting it all together, here is what a routine edit looks like:

\begin{enumerate}
    \item Open the project folder in VS Code.
    \item Open the relevant \file{.qmd} file and start editing.
    \item Run \cmd{quarto preview} (or click the Preview button) to watch
    your changes as you make them.
    \item Once you're happy with the result, stop the preview, run
    \cmd{quarto render} to rebuild the whole site cleanly, and skim a few
    pages to make sure nothing else broke.
    \item Commit and push your changes (Section~\ref{sec:publish}).
    \item Check the live website after a minute or two to confirm your
    change appears correctly.
\end{enumerate}

\section{Editing the Website with Claude Code}
\label{sec:claude-code}

Everything above describes doing the work by hand: opening a file,
editing Markdown and HTML yourself, running the commands yourself. That's
essential to understand, but it isn't the only way to work. \term{Claude
Code} is an AI coding assistant from Anthropic that runs inside VS Code
(or a terminal) and can read the whole project, make the edits itself,
run \cmd{quarto render}, and even commit and push --- all from
plain-English instructions you give it, like ``add a new job listing for
a Research Intern position, same format as the existing ones.''

Starting in Part II, this guide will show most tasks \emph{two} ways: how
to do them yourself, step by step, and a suggested prompt you could give
Claude Code instead to have it done for you in seconds. Understanding the
manual steps is still worthwhile --- it means you can judge whether
Claude Code's changes look right and fix small things yourself when
needed --- but for most day-to-day updates, describing what you want is
often faster than doing it by hand.

\subsection{Installing Claude Code in VS Code}

\begin{enumerate}
    \item Open VS Code and click the Extensions icon in the sidebar (or
    press \cmd{Cmd+Shift+X} / \cmd{Ctrl+Shift+X}).
    \item Search for \term{Claude Code} and install the extension
    published by Anthropic.
    \item A new Claude icon will appear in the sidebar. Click it to open
    the Claude Code panel.
    \item You'll be prompted to sign in. This uses the same kind of
    browser-based sign-in as GitHub in Section~\ref{sec:git-auth}: click
    \textbf{Sign in}, log into your Claude account in the browser window
    that opens, and authorize VS Code.
    \item Using Claude Code requires an active Claude plan (Pro or Max)
    or billing set up on the Anthropic Console. Ask your supervisor if
    ODEKA already has an account you should use, or if you need to set up
    your own at \url{https://claude.ai}.
\end{enumerate}

\subsection{Working with Claude Code}

\begin{enumerate}
    \item Make sure the ODEKA website folder is the one open in VS Code
    (\textbf{File} $\to$ \textbf{Open Folder...} if it isn't) --- Claude
    Code works on whichever project is currently open.
    \item In the Claude Code panel, type what you want in plain English,
    for example:
    \begin{quote}
    \textit{``Add a news post announcing that ODEKA hired a new field
    coordinator, dated today. Use the same format and image size as the
    other news posts.''}
    \end{quote}
    \item Claude Code will make the edits, and can also run
    \cmd{quarto render}/\cmd{quarto preview} for you and describe what it
    changed. Review its summary and the changed files before moving on.
    \item Ask it to show you a preview, or open one yourself
    (Section~\ref{sec:quarto-basics}), to confirm the result looks right.
    \item You can ask Claude Code to commit and push the change too
    (``commit this and push it''), or do that step yourself once you're
    happy with it --- either is fine.
\end{enumerate}

\paragraph{A few habits worth keeping.} Be specific about which page and
what you want (naming the exact page, section, or person makes for a
better result than a vague request). Always look over what changed before
pushing, the same way you'd proofread something before sending it. And if
a result isn't quite right, just say so in plain language --- ``move that
a bit to the left,'' ``make the text smaller'' --- the same way the rest
of this guide describes fixing things by hand.

\subsection{What Claude Code Already Knows: \file{CLAUDE.md}}
\label{sec:claude-md}

A conversation with Claude Code is not remembered afterward. Close the
window, and the next session --- even yours, even five minutes later ---
starts from nothing except the files in the project. That's normally fine,
since Claude Code can read the code itself, but some things aren't visible
just from reading the code once: a bug that only shows up on the live site
and not while previewing locally, a design choice that was tried and
deliberately rejected, a CSS rule that looks redundant but isn't. Losing
that kind of hard-won context between sessions is wasteful --- it means
re-discovering the same problem twice.

To fix this, the repository has a file at its root named \file{CLAUDE.md}.
Claude Code reads it automatically, every time, at the start of any
session opened in this folder --- on any computer, for anyone. It exists
purely to carry exactly this kind of knowledge forward: known traps (for
example, an image that displays fine locally but silently fails once
published, because of a filename case mismatch that only matters on the
live server), the reasoning behind non-obvious choices, and pointers to
this guide for the step-by-step version of a task.

\paragraph{You don't need to write it yourself.} If Claude Code hits
something tricky while helping you --- a bug with a non-obvious cause, a
setting that has to be right in three different files at once --- it
should be the one updating \file{CLAUDE.md}, not you. A reasonable habit:
after Claude Code fixes something confusing, ask it directly, ``is this
worth adding to CLAUDE.md?'' If a future problem turns out to be something
Claude Code already ran into before but never wrote down, that's a sign to
ask it to add a note next time.

\paragraph{How this differs from this guide.} \file{CLAUDE.md} is written
\emph{for Claude Code} --- terse, technical, assumes familiarity with the
code. This document is written \emph{for you} --- the person doing the
work, whether by hand or by asking Claude Code. They're meant to stay in
sync but serve different readers: if you make a structural change to how a
page works, both may need updating (Section~\ref{sec:doc-upkeep} says more
about keeping this document itself current).

\part{The ODEKA Website, Page by Page}

This part walks through every page of the site and the specific edits
you'll actually make: updating a bio, adding a project, publishing a
news post, posting a job, or swapping out a photo. Each task is shown
the manual way first, then --- where it's genuinely faster --- as a
Claude Code prompt (Section~\ref{sec:claude-code}). Reading the manual
version at least once is worth it even if you always use Claude Code
afterward: it's what lets you tell whether what Claude Code did is
actually right.

A general note before starting: after \emph{any} edit described below,
the workflow is always the same three steps from Part I --- render (or
preview while you work), check it, then commit and push
(Section~\ref{sec:publish}). This part won't repeat that every time.

\section{The People Page}
\label{sec:people-page}

The People page (\file{02\_people/people.qmd}) is one long file. It's
organized into named sections (Senior Management, PhD Students, Past
Research Associates, and so on), each containing a run of individual
entries. There's no listing magic here like on News or Jobs --- every
person is hand-written directly in this one file, in the order they
appear on the page.

\subsection{How an Entry Is Built}

Every person has the same two-part shape: a small square photo, and an
info block next to it with their name, title, and (optionally) a short
bio and a link. Most entries look like this:

\begin{lstlisting}
<div style="clear: both;">
</div>
<br>

<div class="person-photo">
  <img src="faces/lastname.jpg" alt="Full Name">
</div>

<div class="person-info">
<div class="person-name">
  Full Name
</div>

<div class="person-title">
  Their Title
</div>

<p>Optional one- or two-sentence bio.

<div class="person-link">
<a href="https://example.com" target="_blank">
  Profile Link Text
</a>
</div>
</div>
\end{lstlisting}

A few sections list people with no photo at all --- currently ``Past
Research Interns'', and ``Research Interns'' should follow the same
pattern once people are added there (it's empty for now). Those use a
slightly different wrapper --- the photo \file{<div>} is simply
skipped, and \file{person-info} gets a second class:

\begin{lstlisting}
<div style="clear: both; margin-bottom: 38px;">
</div>

<div class="person-info person-info--no-photo">
<div class="person-name">
  Full Name
</div>
<div class="person-title">
  Their Title
</div>
<div class="person-link">
<a href="https://example.com" target="_blank">
  LinkedIn Profile
</a>
</div>
</div>
\end{lstlisting}

If someone doesn't have a photo yet, use
\file{faces/placeholder.jpg} rather than leaving it out --- every entry
on the page currently has an image, and the placeholder (a plain
silhouette icon) is what keeps the layout consistent until a real photo
is added.

\subsection{Editing Someone's Information}

Open \file{02\_people/people.qmd}, search for their name, and edit the
text directly --- the name, title, bio paragraph, or link text/URL are
all plain text inside the block shown above. To swap their photo,
replace the file at \file{02\_people/faces/whatever.jpg} with the new
image \emph{using the same filename} (simplest option), or add a new
image file and update the \file{src="..."} path to match.

\begin{quote}
\textit{Ask Claude Code: ``On the People page, update Jon Weigel's
title to `Scientific Director' and change his bio to say he's now at
[institution].''}
\end{quote}

\subsection{Adding a Person}

\begin{enumerate}
    \item Find the section they belong in, and find another entry in
    that same section to use as a template.
    \item Copy that whole entry (from its opening
    \cmd{<div style="clear: both...}} down to the closing
    \cmd{</div>} of \file{person-info}) and paste it either right before
    or right after it, depending on where the new person should sit.
    \item Update the name, title, bio, link, and photo path. If you
    don't have a photo yet, point it at
    \file{faces/placeholder.jpg}.
    \item If the section is ordered a particular way (PhD Students is
    alphabetical by family name; Past Research Associates run
    newest-to-oldest), place the new entry in the position that keeps
    that order.
\end{enumerate}

\begin{quote}
\textit{Ask Claude Code: ``Add a new person to the PhD Students section
of the People page: [Name], PhD Student, no bio yet, placeholder photo.
Keep the section in alphabetical order by family name.''}
\end{quote}

\subsection{Removing a Person}

Delete their entire block, from the \cmd{<div style="clear: both...}}
spacer immediately before their photo down through the closing
\cmd{</div>} of \file{person-info}. Be careful with the closing
\file{</div>} tags in particular --- there are several nested divs in
each entry, and deleting one too many (or too few) will silently break
the layout of the \emph{next} person's entry rather than produce an
error. Render the page afterward and check that both the deleted
person is gone and the people around them still look normal.

\begin{quote}
\textit{Ask Claude Code: ``Remove Jean-François Coté's entry from the
Past Research Associates section of the People page.''}
\end{quote}

\subsection{Adding a New Section}

Copy this pattern, placing it wherever the new section should sit
between two existing ones:

\begin{lstlisting}
<!-- ============================ -->
<!-- ============================ -->
<!-- ============================ -->
<div class="about-section-divider"></div>

# New Section Title {.people-eyebrow-heading}

One-sentence description of who belongs in this section.
<!-- ============================ -->
\end{lstlisting}

Then add person entries underneath it as described above. Two things
to know:

\begin{itemize}
    \item The heading text becomes an entry in the ``on this page'' box
    on the right automatically --- nothing else to configure.
    \item If the new heading's text is identical to another heading
    already on the page (this has come up before, e.g. two sections
    both wanting to be called ``Past Directors''), add an explicit
    anchor to keep their links from colliding: \cmd{\# New Section
    Title \{\#a-unique-id .people-eyebrow-heading\}}.
\end{itemize}

\begin{quote}
\textit{Ask Claude Code: ``Add a new section to the People page called
`Visiting Researchers', right after Scientific Advisory Committee, with
the description `ODEKA hosts visiting researchers who contribute to our
work during their stay.' Leave it empty for now.''}
\end{quote}

\section{The Projects Page}
\label{sec:projects-page}

Unlike People, Projects uses Quarto's \term{listing} system: each
project is its own file in \file{03\_projects/posts/}, and both the
Projects page and the homepage carousel are automatically generated
from those files based on a few frontmatter fields. You never edit a
list of projects by hand --- you add, edit, or delete the individual
project files, and the listings update themselves.

\subsection{The Fields That Matter}

Every project file's frontmatter looks like this:

\begin{lstlisting}
---
title: "Project Name"
author:
  - Person One
  - Person Two
date: 2020-01-01
display date: "2020-2023"
image: ../images/projectN.jpg
description: "One-sentence summary shown on the card."
status: ongoing
type: major
about:
  template: marquee
---
<style>
header.quarto-title-block .description,
.quarto-title .description,
.quarto-title-meta .description,
p.description {
  display: none !important;
}
</style>

[← Back to Projects](../index.qmd){.breadcrumb}

The actual page content goes here.
\end{lstlisting}

The \file{marquee} template (the same one News posts use) is what puts
a large banner image at the top of the page instead of a small sidebar
photo --- the \cmd{<style>} block right after it is required for the
same reason it's required on News posts (Section~\ref{sec:news-page}):
without it, the \file{description} field shows up twice.

Two fields drive everything about where a project shows up, and they
are independent of each other:

\begin{itemize}
    \item \term{\file{type}}: either \cmd{major} or \cmd{other}. This
    decides which section of the \emph{Projects page itself} the
    project appears under (``Major Projects'' or ``Other Projects'').
    \item \term{\file{status}}: currently \cmd{ongoing} or anything
    else (e.g. \cmd{completed}). This only affects the little
    ``ONGOING'' badge shown on the card --- \emph{except} for one
    important consequence below.
\end{itemize}

\subsection{The Homepage Carousel Is Pickier Than the Projects Page}

The Projects page shows every project of a given \file{type}, regardless
of status. The \term{homepage} carousel (``Ongoing major projects'') is
more restrictive: it only shows projects where \file{type: major}
\emph{and} \file{status: ongoing} are \emph{both} true. This means:

\begin{itemize}
    \item Marking a major project's status as anything other than
    \cmd{ongoing} removes it from the homepage immediately, but it
    stays exactly where it was on the Projects page.
    \item An \cmd{other}-type project will never appear on the homepage
    carousel, no matter its status.
    \item If you want a project to stop being featured on the homepage
    the moment it wraps up, changing its \file{status} field is the
    single edit that does it --- nothing else needs to change.
\end{itemize}

\subsection{Adding a Project}

\begin{enumerate}
    \item Create a new file in \file{03\_projects/posts/}, e.g.
    \file{my-new-project.qmd}, using the template above.
    \item Add a card image to \file{03\_projects/images/} and point
    \file{image:} at it.
    \item Set \file{type} and \file{status} deliberately, per the rules
    above.
    \item Write the actual project page content below the frontmatter
    (see Section~\ref{sec:projects-todo} for the current state of this
    --- it's genuinely unresolved for existing projects too).
\end{enumerate}

\begin{quote}
\textit{Ask Claude Code: ``Add a new major, ongoing project called
`[Name]' to the Projects page, described as `[one sentence]'. Use
[image file] as the card image.''}
\end{quote}

\subsection{Editing or Removing a Project}

To edit, open its file in \file{03\_projects/posts/} and change the
frontmatter fields or the page content directly. To remove a project
entirely, delete its \file{.qmd} file (and its image, if nothing else
uses it) --- it disappears from every listing automatically.

\begin{quote}
\textit{Ask Claude Code: ``Change the Religion \& Development project's
status to completed'' or ``Delete the [project name] project.''}
\end{quote}

\section{The News Page}
\label{sec:news-page}

News works the same way as Projects: one file per post in
\file{04\_news/posts/}, listed automatically.

\subsection{Adding, Editing, or Removing a Post}

A news post's frontmatter is simpler than a project's:

\begin{lstlisting}
---
title: "Headline Goes Here"
author:
  - Author Name
date: "Month DD, YYYY"
description: "One-sentence summary."
image: ../images/newsYYYYMMDD.jpg

about:
  template: marquee
---
<style>
header.quarto-title-block .description,
.quarto-title .description,
.quarto-title-meta .description,
p.description {
  display: none !important;
}
</style>

[← Back to News](../index.qmd){.breadcrumb}

**A bold one-sentence lead-in, restating the summary.**

The rest of the story, in normal paragraphs.
\end{lstlisting}

That \cmd{<style>} block isn't decorative --- without it, the
\file{description} shows up twice (once as Quarto's own subtitle, once
as your bold lead-in paragraph). Keep it in every new post.

To add a post: copy an existing file in \file{04\_news/posts/}, rename
it (the convention so far is \file{newsYYYYMMDD.qmd}), update the
frontmatter and body, and add its image to \file{04\_news/images/}. To
edit one, just open and change it. To remove one, delete its file.

\begin{quote}
\textit{Ask Claude Code: ``Add a news post dated today announcing
[whatever happened]. Use the same format as the existing posts.''}
\end{quote}

\subsection{Turning the News Section On and Off}
\label{sec:news-switch}

If nobody is actively writing news, the whole section can be switched
off --- hidden from the navbar, removed from the homepage (replaced by
a plain divider), and excluded from the built site --- without deleting
any of the actual posts. Turning it back on later restores everything
exactly as it was, because nothing in \file{04\_news/} itself is ever
touched by the switch. It only takes a handful of small, coordinated
edits across two files:

\begin{enumerate}
    \item In \file{\_quarto.yml}, find the block of comments marked
    \cmd{NEWS SECTION SWITCH} near the top, inside the project's
    \file{render:} list. Uncomment the single \cmd{"!04\_news/**"}
    exclusion line inside it (the list already always excludes
    \file{CLAUDE.md} above it --- leave that line alone).
    \item A little further down in the same file, in
    \file{navbar: left:}, comment out the two lines for the News entry
    (marked with a \cmd{NEWS SWITCH} comment right above them).
    \item In \file{index.qmd}, two more places are marked
    \cmd{NEWS SWITCH}:
    \begin{itemize}
        \item Comment out the \file{news\_carousel} listing definition
        in the frontmatter (this one matters --- if you skip it, Quarto
        will dump the news carousel awkwardly at the very bottom of the
        homepage instead of hiding it).
        \item Change the \cmd{\{\{< include ... >\}\}} line just above
        where ``Latest News'' used to be from
        \file{news\_section\_on.qmd} to \file{news\_section\_off.qmd}.
    \end{itemize}
    \item Run \cmd{quarto render}. That single render does everything:
    removes the News pages that were already built, updates the
    navbar, and swaps in the plain divider on the homepage.
\end{enumerate}

To turn News back on, reverse the edits from steps 1--3 above
(re-comment/uncomment the opposite way), then render again.

\begin{quote}
\textit{Ask Claude Code: ``Turn the News section off'' or ``Turn the
News section back on'' --- it knows to make all the coordinated edits
together and re-render.}
\end{quote}

See Section~\ref{sec:news-todo} for a note on \emph{when} you might
actually want to do this.

\section{The Jobs Page}
\label{sec:jobs-page}

Jobs also uses one file per posting, in \file{05\_jobs/posts/}, but the
page layout is richer than News --- it has a structured header
(supervisors, dates, location, duration) above the actual posting text.

\subsection{The Content Rule: Copy the Posting Verbatim}

This is important enough to call out on its own: when you add a real
job posting, the body text (Overview, Responsibilities, Qualifications,
etc.) should be an exact, word-for-word copy of the actual posting
--- wherever it's really hosted (J-PAL/IPA, a university jobs board,
wherever) --- not a rewritten or summarized version. Only the
\emph{layout} (the section headers, the meta block at the top, colors,
fonts) is ours; the wording is not. When you fetch the source page to
copy from, copy the literal text, not a paraphrase of it.

One consequence of this: our postings are sometimes hosted on a
recruiting platform (J-PAL/IPA's shared application system) as a matter
of practicality, but \term{ODEKA is always the actual employer}. Keep
that clear on the page --- e.g. the short description under the title
should say ``We are hiring...'', not name the hosting platform as the
employer.

\subsection{The Template}

\begin{lstlisting}
---
title: "Exact Job Title"
author:
  - Supervisor Name, Affiliation
date: "Month DD, YYYY"
description: "We are hiring a [role] to [one-sentence summary]."
language:
  title-block-author-single: "Supervisor"
  title-block-author-plural: "Supervisors"
  listing-page-field-author: "Supervisor"
---

<div class="job-page">

<div class="job-meta">
<div class="job-meta-item"><strong>Posted</strong>Month DD, YYYY</div>
<div class="job-meta-item"><strong>Supervisors</strong>Name One<br>Name Two</div>
<div class="job-meta-item"><strong>Location</strong>City, Country</div>
<div class="job-meta-item"><strong>Start Date</strong>Month D, YYYY</div>
<div class="job-meta-item"><strong>Duration</strong>N months</div>
</div>

[← View All Open Positions](../index.qmd){.breadcrumb}

### Overview
Copied verbatim from the real posting.

### Responsibilities
- Copied verbatim, as a bullet list.

### Qualifications
- Copied verbatim (required)
- Copied verbatim (desired)

### Supervisors
**Primary supervisors:** ...
**Secondary supervisors:** ...

### How to Apply
Copied verbatim, plus a plain-language note that ODEKA is the employer
if the original posting doesn't make that obvious.

<div class="job-apply">
<a class="about-button" href="https://actual-application-url"
   target="_blank">Apply for this position →</a>
</div>

</div>
\end{lstlisting}

The \file{language:} block is what relabels Quarto's default ``Author''
heading to ``Supervisor''/``Supervisors'' --- keep it in every posting.
The section headers (\cmd{\#\#\#}) can be adjusted to match whatever
structure the real posting actually uses; they don't have to be exactly
these five.

\subsection{Adding, Editing, or Removing a Posting}

Same pattern as News and Projects: add a new file in
\file{05\_jobs/posts/} for a new posting, edit a file directly to
update one, delete a file to remove one (once it's filled or closed).

\begin{quote}
\textit{Ask Claude Code: ``Here's a job posting URL: [link]. Add it to
the Jobs page, copying the content verbatim, in our usual job-page
format.''}
\end{quote}

\section{Contacts and Legal Notice}

These two are the simplest pages on the site --- each is a single file
(\file{06\_contacts/contacts.qmd} and
\file{07\_legal\_notice/legal\_notice.qmd}) with plain text content,
organized under \cmd{\#\#\#} headings:

\begin{lstlisting}
### Section Heading:
Whatever text or contact details go here.<br>
Use <br> for a line break within a section.
\end{lstlisting}

Just open the relevant file and edit the text directly --- add, remove,
or reword a \cmd{\#\#\#} section however needed. The photo shown next to
the text is set by the \file{image:} field in the frontmatter
(\file{about: image: image1.jpg}); replace that file to change it.

\section{Changing an Illustration Anywhere on the Site}

Every image on the site is a plain file referenced by path, either in a
page's frontmatter (\file{image: ...}) or directly in the page content
(\file{<img src="...">}). To swap one out, you have two options:

\begin{enumerate}
    \item \term{Simplest}: replace the file at that exact path with
    your new image, keeping the same filename. Nothing else needs to
    change.
    \item Add a new image file and update the path that references it,
    if you'd rather keep the old one around or use a more descriptive
    filename.
\end{enumerate}

\subsection{Where Images Live, and How They're Named}

Every page keeps its own images in its own folder, right next to the
\file{.qmd} file that uses them --- there is no single shared
\file{images/} folder for the whole site. The pattern is consistent
enough to state as a table:

\begin{center}
\begin{tabular}{ll}
\toprule
\textbf{What} & \textbf{Where, and how it's named} \\
\midrule
Page banner (top/bottom of a page) & \file{0N\_folder/imageN.ext}
(sequential) \\
Person photo & \file{02\_people/faces/lastname.ext} \\
Project card image & \file{03\_projects/images/projectN.ext} \\
News post image & \file{04\_news/images/newsYYYYMMDD.ext} \\
Site-wide logos/brackets & root folder, \file{odeka\_logo\_*.png} \\
\bottomrule
\end{tabular}
\end{center}

\term{Stick to whichever pattern already applies to the page you're
editing.} This matters more than it might seem, for two reasons we've
already hit in practice, not hypothetically:

\begin{itemize}
    \item \term{File extension case has to match exactly.} One news
    post once linked to \file{news20220627.JPG} while the actual file
    on disk was \file{news20220627.jpg}. It worked by accident on a Mac
    (whose filesystem ignores case) and would have silently shown a
    broken image once deployed, since GitHub Pages' hosting does not.
    Always match the real filename's case exactly, and prefer lowercase
    extensions for anything new.
    \item \term{Unused images pile up fast if filenames drift.} A
    routine cleanup pass on this site once found --- and removed --- ten
    image files that nothing on the site referenced anymore, plus an
    entire folder of raw, never-used originals. That happens naturally
    when a new image doesn't follow the existing naming pattern closely
    enough to make it obvious which page it belongs to, or whether an
    old one it was meant to replace is still hanging around.
\end{itemize}

Following the table above --- same folder as the page, same naming
style as the images already there --- is what keeps that from
happening again. If you ever do add an image and later decide not to
use it, delete it rather than leaving it behind.

\subsection{Cropping and Sizing Conventions}
\label{sec:cropping}

\begin{itemize}
    \item \term{Person photos} (\file{02\_people/faces/}): square
    crops, centered on the face, saved at a reasonable resolution (the
    existing photos are roughly 400--800px square). They're displayed
    at 200×200px, so anything comfortably larger than that is fine.
    \item \term{Banners} (the full-width images at the top and bottom
    of most pages): wide landscape crops. Keep the subject roughly
    centered horizontally --- the banner crops responsively as the
    browser window resizes, and content too close to the left or right
    edge can get cut off on narrower screens.
    \item \term{Card images} (Projects, News): roughly a 4:3 or similar
    landscape crop works best; they display inside a fixed-size card.
\end{itemize}

\begin{quote}
\textit{Ask Claude Code: ``Replace [person]'s photo on the People page
with [new file], cropped the same way as the others.''}
\end{quote}

\section{What Remains To Be Done}

This section is a snapshot of known gaps as of this writing --- a
starting checklist, not a complete one. Update or trim it as items get
resolved.

\subsection{People Page}

\begin{itemize}
    \item \term{Full names missing}: several Senior Management and
    Senior Enumerators entries currently show only a first name
    (Junior, Gilbert, Jacques; Max, Grégoire, Bernard, Serge, André,
    Benjamin, Théo, Jean). Family names need to be added.
    \item \term{Photos missing}: the people listed above, plus most of
    the current PhD Students (Clotaire Boyer, Sarah Danner, Eva
    Davoine, Gabriel Granato, Nic Wicaksono), are still showing the
    placeholder silhouette. Arthur Laroche, the sixth current PhD
    student, already has a real photo --- his entry is the model to
    follow for the rest. See Section~\ref{sec:cropping} for cropping
    conventions once more real photos are ready.
    \item \term{PhD Students section is missing information}: for each
    current PhD student, we're missing the year they started working
    with ODEKA and their PhD institution --- probably worth labeling
    the same way Past PhD Students already does it (a \cmd{**PhD:**}
    line), but that's worth double-checking once the actual information
    is in hand rather than assuming the format transfers directly.
    \item \term{Research Associates and Research Interns sections are
    empty}: both currently exist as headings with a description and no
    one listed underneath. Add current members as they're identified
    --- Research Interns should use the no-photo entry pattern
    (Section~\ref{sec:people-page}), matching Past Research Interns.
    \item \term{Past People}: double-check whether anyone who has left
    ODEKA is missing from the appropriate ``Past...'' section.
\end{itemize}

\subsection{Projects Page}
\label{sec:projects-todo}

Every project's actual page content is currently a placeholder
(``Page under construction''). Beyond just writing that content, the
following are genuinely open questions --- not yet decided, and worth
real thought rather than a quick default:

\begin{itemize}
    \item What a project page should actually contain, and how it
    should be laid out.
    \item What the illustration strategy should be (one image per
    project? several? a gallery?).
    \item How to choose what image represents a given project.
\end{itemize}

One specific open question, unresolved as of this writing: Jon has
asked for articles to be listed \emph{by theme} somewhere on the site,
but the current structure is organized by \emph{project}, and a theme
is not the same thing as a project (a single theme could span several
projects, or a project could touch several themes). How --- or whether
--- to reconcile ``browse by project'' with ``browse by theme'' hasn't
been figured out yet.

It's worth being precise about what ``project'' means here in the
first place: from ODEKA's side, what we've been calling a project is
really an RCT program --- a specific randomized evaluation, not a
research topic. Organizing the site around that unit seems like the
intuitive choice for what is, at bottom, a program-evaluation research
platform: it's the natural unit of \emph{our} work. But that intuition
is exactly the kind worth distrusting a little --- what's obvious from
ODEKA's side of the platform isn't necessarily what's most useful to
the researchers actually browsing it, and their sense of what should
group together (by theme, by question, by region, by something else
entirely) may not line up with the program structure at all. This is a
real design question, not just a missing feature, and deserves to be
thought through with that in mind rather than defaulted one way or the
other.

\subsection{News Page}
\label{sec:news-todo}

The intended model for this section is that people out in the field
write their own news posts and sign them in their own name --- a
first-person record of their time at ODEKA, not just an announcements
feed. If someone in the field wants to do that, they should keep
writing.

Getting credit for what you wrote, under your own name, is itself part
of the motivation for anyone to want to keep this up --- so the byline
matters, not just the content. As of this writing, the existing posts
are attributed to ``The Research Associate Team'' rather than the
individual who actually wrote them (Adrien, for all of them). If the
sign-your-own-name model above is the one that's actually adopted
going forward, those existing posts' \file{author:} fields should be
updated to match --- back to the person who wrote each one --- rather
than left as a team byline while new posts are signed individually
(Section~\ref{sec:news-page}).

If nobody is willing to keep the section updated, though, that's a
real, acceptable option too: rather than let it visibly go stale,
turn it off (Section~\ref{sec:news-switch}). Nothing is lost by doing
that --- every post that's already been written stays intact and can be
switched back on the moment someone wants to pick it back up.

\subsection{Custom Domain, and Where the Address Should Land}
\label{sec:domain-todo}

The site currently lives at its default GitHub Pages address
(\url{https://odeka-asbl.github.io}), not a domain ODEKA owns. Two
separate things need deciding here, and they shouldn't be conflated:

\begin{itemize}
    \item \term{Linking a custom domain.} This is mostly mechanical once
    a domain is chosen and registered: GitHub Pages supports it via a
    \file{CNAME} file naming the domain, plus a DNS record set up
    wherever the domain was purchased (either a \cmd{CNAME} record
    pointing at \file{odeka-asbl.github.io}, or \cmd{A} records
    pointing at GitHub's IP addresses, depending on whether it's a
    subdomain or an apex/root domain). The \file{CNAME} file itself
    should live in the \emph{project root} (alongside
    \file{\_quarto.yml}), not be placed by hand inside \file{docs/} ---
    Quarto automatically copies loose files like this from the project
    root into \file{docs/} on every render (the same way it already
    does for \file{favicon.ico}), so a root-level copy survives
    renders reliably where a \file{docs/}-only one wouldn't.
    Registering the domain itself and configuring DNS happens outside
    this repository, with whoever manages ODEKA's domain registrar ---
    Claude Code can help with the \file{CNAME} file and the DNS record
    values once a domain is picked, but can't register a domain or log
    into a registrar on its own.
    \item \term{Where the address should land.} Separately, and still
    undecided: Jon has asked that the website address land directly on
    the \emph{Projects} page rather than the current homepage. That's
    not automatic just from pointing a domain at the site --- as
    things stand, the root address serves \file{index.qmd} (the
    homepage), and Projects is a page underneath it. Making Projects
    the landing page means an actual decision about \emph{how}: whether
    the homepage should redirect to Projects, whether Projects should
    effectively \emph{become} the homepage (absorbing whatever from the
    current homepage is still worth keeping --- the project carousels,
    the quotes, the News section), or whether the navbar and homepage
    both stay as they are and only the bare domain root redirects. This
    is a real design question in the same spirit as the
    project-versus-theme one above (Section~\ref{sec:projects-todo}),
    not a small technical toggle --- worth resolving deliberately once
    a domain is actually in hand, rather than picking the easiest
    option by default.
\end{itemize}

\subsection{Keeping This Document Itself Up to Date}
\label{sec:doc-upkeep}

As items on this list get done --- full names added, photos in place,
a real project page written --- this guide should be updated to match,
not left describing a state of the site that no longer exists. That's
exactly the kind of small, well-defined edit Claude Code is good at:
point it at whatever changed and ask it to update the relevant part of
\file{read\_me/odeka\_website\_guide.tex} accordingly (it can even
recompile the PDF for you).

The same goes for \file{CLAUDE.md} (Section~\ref{sec:claude-md}): the two
documents are meant to stay in sync. As a rough rule, if a fix or a
decision would help a future \emph{person} doing this work, it belongs
here; if it would help a future \emph{Claude Code session} avoid
re-debugging the same problem, it belongs in \file{CLAUDE.md}. Often it's
worth adding to both.

\end{document}
